% NeurIPS-style (or ICML-style) 8-page paper skeleton
% ---------------------------------------------------
% NOTE: For NeurIPS, you typically use the neurips_YYYY style.
% Replace YYYY with the target year once you know it (e.g., neurips_2024).
% This file is a clean starting point that compiles with standard NeurIPS style.

\documentclass{article}

% If you have the NeurIPS style file:
% \usepackage[final]{neurips_2024} % or [preprint], [final]
% Otherwise this falls back to plain article formatting.
\IfFileExists{neurips_2024.sty}{
  \usepackage[preprint]{neurips_2024}
}{
  \usepackage[margin=1in]{geometry}
}

% --------------------
% Packages
% --------------------
\usepackage{amsmath,amssymb,amsfonts}
\usepackage{graphicx}
\usepackage{booktabs}
\IfFileExists{microtype.sty}{\usepackage{microtype}}{}
\usepackage{hyperref}
\usepackage{url}
\usepackage{xcolor}
\usepackage[ruled,vlined,linesnumbered]{algorithm2e}

% Theorems
\usepackage{amsthm}
\newtheorem{theorem}{Theorem}
\newtheorem{lemma}{Lemma}
\newtheorem{proposition}{Proposition}
\newtheorem{corollary}{Corollary}
\theoremstyle{definition}
\newtheorem{definition}{Definition}
\newtheorem{remark}{Remark}

% --------------------
% Macros
% --------------------
\newcommand{\R}{\mathbb{R}}
\newcommand{\E}{\mathbb{E}}
\newcommand{\KL}{\mathrm{KL}}
\newcommand{\softmax}{\mathrm{softmax}}
\newcommand{\indic}{\mathbb{I}}
\newcommand{\Dom}{\mathrm{Dom}}
\newcommand{\tr}{\mathrm{tr}}
\providecommand{\coloneqq}{:=}

% --------------------
% Title
% --------------------
\title{Variational Partial Orders for Traces and Structure Learning}

\author{
Anonymous authors\\
Paper under double-blind review
}

\begin{document}

\maketitle

\begin{abstract}
Partial orders arise in two ubiquitous regimes: (i) \emph{execution traces}, where a latent
poset is observed through noisy linearizations, and (ii) \emph{graph structure learning}, where
domain experts provide partial-order priors over variables. These regimes have developed largely
separate toolkits: probabilistic poset generative models with expensive inference, and differentiable
DAG learning with partial-order constraints that can be computationally prohibitive when enforced
na{\"i}vely via transitive closure. We propose a unified variational perspective: (1) a latent-embedding
poset model for de-linearizing traces with a frontier-restricted Plackett--Luce likelihood, (2) an
efficient representation of partial-order priors via transitive reduction and maximal-path factoring to
avoid quadratic constraint blow-up, and (3) scalable variational inference objectives and gradients,
including importance-weighted VI with REINFORCE/VIMCO-style estimators when reparameterization
is unavailable. The result is a plug-and-play route to amortized, uncertainty-aware poset inference
and prior-constrained structure learning across domains.
\end{abstract}

\section{Introduction}
Partial orders (posets) provide a compact representation of precedence constraints and concurrency.
They appear in \emph{trace de-linearization} problems (workflow logs, tool-augmented LLM traces,
robotic assembly), where observed sequences are noisy linear extensions of a latent partial order; and
in \emph{structure learning} problems (causal discovery), where experts often know partial order
constraints (e.g., gene activation cascades, treatment protocols, seasonal patterns).

Despite this shared object, current methods split across regimes. Trace models typically treat the poset
as a latent discrete object and rely on MCMC, while differentiable DAG learning optimizes continuous
graph parameters in the DAG space and historically lacked efficient ways to incorporate general
partial-order priors. Recent work shows that enforcing partial orders via path prohibitions over the
transitive closure can incur $\Theta(m^2)$ terms for a chain of length $m$, making long ordering priors
impractical, and proposes an augmented-acyclicity alternative based on transitive reduction and maximal
paths.

This paper draft provides a unified NeurIPS-style narrative:
\begin{itemize}
\item We formalize a latent-embedding poset model for traces (your HPO) and derive variational objectives.
\item We show how to incorporate partial-order priors efficiently (avoid $\Theta(|O^+|)$ penalties) using maximal-path factoring.
\item We present VI algorithms for this setting, including non-reparameterized importance-weighted VI with VIMCO-style gradients.
\end{itemize}

\section{Background and Problem Setup}

\subsection{Partial orders, transitive closure, and transitive reduction}
Let $O \subseteq [m]\times[m]$ encode partial-order relations $i \succ j$ among $m$ actions/variables.
Its transitive closure $O^+$ contains all implied relations (if $i\succ j$ and $j\succ k$ then $i\succ k$).
The transitive reduction $O^-$ is the smallest relation satisfying $(O^-)^+ = O^+$.

\begin{remark}[Quadratic blow-up for long chains]
For a single chain $v_1 \succ v_2 \succ \cdots \succ v_m$, the transitive closure contains all pairs
$(v_a,v_b)$ with $a<b$, hence $|O^+| = \binom{m}{2} = \Theta(m^2)$. Any constraint or penalty that
enumerates all $(i,j)\in O^+$ inherits this quadratic scaling.
\end{remark}

\subsection{Two regimes: traces vs.\ IID samples}
\textbf{Trace regime:} We observe sequences $y^{(i)}$ that are noisy linear extensions of a latent poset $h$.
\textbf{Structure-learning regime:} We observe samples $D$ from an SEM/Bayesian network and seek a DAG
(or weighted adjacency $W$) while respecting an ordering prior $O$.

\section{Latent Posets from Traces (HPO Model)}

\subsection{Prior over latent embeddings and induced poset}
Let $m$ be the number of actions and $K$ the latent dimension. We place a Gaussian prior over
action embeddings:
\begin{equation}
U_j \sim \mathcal{N}(0, \Sigma_\rho), \quad j=1,\ldots,m,
\label{eq:U_prior}
\end{equation}
where $\Sigma_\rho$ is exchangeable, e.g.\ $\Sigma_\rho=(1-\rho)I+\rho \mathbf{1}\mathbf{1}^\top$, $\rho\in[0,1]$.
We deterministically map $U$ to a strict partial order:
\begin{equation}
h = \Dom(U).
\label{eq:Dom}
\end{equation}

\begin{definition}[Coordinate-wise dominance map]
A simple instantiation is: $i \succ_h j$ iff $U_{i,k} > U_{j,k}$ for all $k\in[K]$.
\end{definition}

\begin{theorem}[Dominance induces a strict partial order]
If $i \succ_h j \iff \forall k,\ U_{i,k} > U_{j,k}$, then $\succ_h$ is irreflexive and transitive, hence $h$ is a strict partial order.
\end{theorem}
\begin{proof}[Proof sketch]
Irreflexive: $U_{i,k} > U_{i,k}$ is impossible. Transitive: coordinate-wise $>$ is transitive.
\end{proof}

\subsection{Frontier-restricted Plackett--Luce likelihood}
Given a trace $y=(y_1,\ldots,y_T)$ and prefix $y_{<t}=(y_1,\ldots,y_{t-1})$, define the remaining actions:
\[
R_t \coloneqq A \setminus \{y_1,\ldots,y_{t-1}\}.
\]
The feasible set (frontier) is the set of minimal elements of $R_t$ under $\succ_h$:
\begin{equation}
F_t(h;y_{<t}) \coloneqq \{a\in R_t: \nexists b\in R_t \text{ such that } b \succ_h a\}.
\label{eq:frontier}
\end{equation}
We model the next action via a frontier-restricted Plackett--Luce / Boltzmann policy with trembling-hand noise:
\begin{equation}
p(y_t \mid y_{<t}, h, \beta, \epsilon) =
(1-\epsilon)\frac{\exp(\beta Q(y_t;h,t))}{\sum_{a\in F_t(h)}\exp(\beta Q(a;h,t))}
+ \epsilon\frac{1}{|R_t|},
\label{eq:frontier_softmax}
\end{equation}
where the first term is treated as $0$ if $y_t\notin F_t(h)$.

A full trace has likelihood:
\begin{equation}
p(y \mid h,\beta,\epsilon) = \prod_{t=1}^T p(y_t \mid y_{<t}, h, \beta,\epsilon).
\label{eq:trace_likelihood}
\end{equation}

\subsection{Posterior}
Let $D=\{y^{(i)}\}_{i=1}^n$ be traces. Unknowns are $(U,\rho,\beta,K)$ (optionally infer $K$). The posterior:
\begin{equation}
p(U,\rho,\beta,K \mid D) \propto p(U\mid\rho,K)p(\rho)p(\beta)p(K)\prod_{i=1}^n p(y^{(i)}\mid h(U),\beta,\epsilon).
\label{eq:posterior}
\end{equation}

\section{Variational Inference for HPO (HPO-VI)}

\subsection{ELBO}
Choose a variational family $q_\phi(U,\rho,\beta,K)$ (e.g.\ mean-field). The ELBO is:
\begin{align}
\mathcal{L}(\phi)
&= \E_{q_\phi}\Big[
\sum_{i=1}^n \log p(y^{(i)}\mid h(U),\beta,\epsilon)
+ \log p(U\mid\rho,K) + \log p(\rho)+\log p(\beta)+\log p(K)
- \log q_\phi(U,\rho,\beta,K)
\Big].
\label{eq:elbo}
\end{align}
The key difficulty is that $h(U)=\Dom(U)$ is piecewise-constant in $U$: within any
connected region of the embedding space the induced poset is fixed, and the function is
discontinuous at boundaries.  Pathwise/reparameterized gradients through the likelihood are
therefore zero almost everywhere, making standard reparameterization tricks ineffective.
We therefore require score-function estimators and importance-weighted objectives.

\subsection{Importance-weighted objectives (VR-IWAE)}
\label{sec:vr_iwae}
Let $z=(U,\rho,\beta,K)$ and draw $z_{1:N}\stackrel{iid}{\sim} q_\phi$.
Define the unnormalized importance weight
\begin{equation}
w(z) \coloneqq \frac{p(D,z)}{q_\phi(z)}
= \exp\!\big(\log p(D,z)-\log q_\phi(z)\big),
\label{eq:imp_weight}
\end{equation}
where
\begin{equation}
\log p(D,z) =
\log p(U\mid\rho,K)+\log p(\rho)+\log p(\beta)+\log p(K)
+ \sum_{i=1}^n \log p(y^{(i)}\mid h(U),\beta,\epsilon).
\label{eq:log_joint}
\end{equation}
Following \citet{burda2016iwae} and \citet{daudel2026importance},
the VR-IWAE family of objectives parameterized by $\alpha\in[0,1)$ is:
\begin{equation}
\mathcal{L}^{(\alpha)}_N(\phi)
= \frac{1}{1-\alpha}\,
\E\!\left[\log\!\left(\frac{1}{N}\sum_{n=1}^N w(z_n)^{1-\alpha}\right)\right].
\label{eq:vr_iwae}
\end{equation}
When $\alpha=0$ this recovers the standard IWAE bound~\citep{burda2016iwae};
for $\alpha>0$ one obtains R\'{e}nyi-divergence bounds analyzed in~\citet{daudel2026importance}.

\subsection{Score-function gradient decomposition}
\label{sec:reinforce_grad}
Define the self-normalized weights
\begin{equation}
W_n^{(1-\alpha)}
\coloneqq
\frac{w(z_n)^{1-\alpha}}{\sum_{j=1}^N w(z_j)^{1-\alpha}}.
\label{eq:snw}
\end{equation}
For any variational parameter block $\psi\subseteq\phi$,
the (exact) score-function gradient of~\eqref{eq:vr_iwae} decomposes as~\citep{daudel2026importance}:
\begin{align}
\nabla_\psi \mathcal{L}^{(\alpha)}_N
&= \E_{z_{1:N}\sim q_\phi}\Bigg[
\underbrace{\sum_{n=1}^N W_n^{(1-\alpha)}\,\nabla_\psi \log w(z_n)}_{\text{self-normalized IS term}}
+ \underbrace{\frac{1}{1-\alpha}\left(\sum_{n=1}^N \nabla_\psi \log q_\phi(z_n)\right)
\log\!\left(\frac{1}{N}\sum_{j=1}^N w(z_j)^{1-\alpha}\right)}_{\text{REINFORCE term (high variance)}}
\Bigg].
\label{eq:reinforce_grad}
\end{align}
Since $p(D,z)$ does not depend on variational parameters $\psi\subseteq\phi$,
we have $\nabla_\psi\log w(z)=-\nabla_\psi\log q_\phi(z)$.
The first term is well-behaved, but the second (REINFORCE) term can exhibit
SNR collapse as $N$ grows~\citep{daudel2026importance,mnih2016vimco},
motivating control-variate methods.

\subsection{VIMCO and VIMCO-$\star$ control variates}
\label{sec:vimco}

\paragraph{General VIMCO baselines.}
The VIMCO framework~\citep{mnih2016vimco} replaces the global log-sum REINFORCE signal
with per-sample ``leave-one-out'' baselines $f^{(\alpha)}_{-n}$ that depend only on $z_{-n}$:
since $\E[\nabla_\psi\log q_\phi(z_n)\,f(z_{-n})]=0$, unbiasedness is preserved.
\citet{daudel2026importance} provide a systematic family of such baselines for general $\alpha$.

\paragraph{VIMCO-$\star$ for IWAE ($\alpha=0$).}
For the standard IWAE case ($\alpha=0$), define $W_n=\frac{w(z_n)}{\sum_{j=1}^N w(z_j)}$.
The VIMCO-$\star$ estimator~\citep{daudel2026importance} takes a particularly simple form:
\begin{equation}
\boxed{
\widehat{\nabla_\psi \mathcal{L}^{(0)}_N}
= \sum_{n=1}^N\Big[
W_n\,\nabla_\psi \log w(z_n)
- \log(1-W_n)\,\nabla_\psi \log q_\phi(z_n)
\Big].
}
\label{eq:vimco_star}
\end{equation}
For variational parameters $\psi\subseteq\phi$ (where $\nabla_\psi\log w=-\nabla_\psi\log q$),
this simplifies to:
\begin{equation}
\widehat{\nabla_\psi \mathcal{L}^{(0)}_N}
= -\sum_{n=1}^N \big(W_n+\log(1-W_n)\big)\,\nabla_\psi \log q_\phi(z_n).
\label{eq:vimco_star_simplified}
\end{equation}

\begin{remark}[SNR scaling]
\citet{daudel2026importance} prove that the signal-to-noise ratio (SNR) of the
VIMCO-$\star$ estimator at $\alpha=0$ scales as $\sqrt{N}$, avoiding the SNR collapse
that afflicts naive multi-sample REINFORCE estimators.
This makes it the recommended default for HPO-VI.
\end{remark}

\begin{remark}[Implementation convenience]
Equations~\eqref{eq:vimco_star}--\eqref{eq:vimco_star_simplified} never require gradients
through $h(U)=\Dom(U)$; only $\log q_\phi(z_n)$ and its gradient w.r.t.\ variational
parameters are needed.  This is the key practical advantage for HPO-VI.
\end{remark}

\subsection{Algorithm}

\begin{algorithm}[t]
\caption{HPO-VI with IWAE ($\alpha=0$) and VIMCO-$\star$ gradients}
\label{alg:hpo_vi}
\KwIn{Traces $D=\{y^{(i)}\}_{i=1}^n$, particles $N$, step size $\eta$, slip $\epsilon$}
Initialize variational parameters $\phi$ for $q_\phi(z)$, $z=(U,\rho,\beta,K)$\;
\While{not converged}{
  Sample $z_{1:N}\sim q_\phi(z)$\;
  \For{$n=1$ \KwTo $N$}{
    $h^{(n)}\leftarrow \Dom(U^{(n)})$\;
    $\log p_n \leftarrow \log p(z_n)+\sum_{i=1}^n \log p(y^{(i)}\mid h^{(n)},\beta^{(n)},\epsilon)$\;
    $\log q_n \leftarrow \log q_\phi(z_n)$\;
    $\log w_n \leftarrow \log p_n - \log q_n$\;
  }
  $W_n \leftarrow \exp(\log w_n)\big/\sum_{j=1}^N \exp(\log w_j)$
  \tcp*[f]{normalize in log-space for stability}\;
  $\widehat{\nabla_\phi}
  \leftarrow -\sum_{n=1}^N\big(W_n+\log(1-W_n)\big)\,\nabla_\phi \log q_\phi(z_n)$
  \tcp*[f]{Eq.~\eqref{eq:vimco_star_simplified}}\;
  $\phi \leftarrow \phi + \eta\,\widehat{\nabla_\phi}$\;
}
\KwOut{Variational posterior $q_\phi(z)$ and samples of $h=\Dom(U)$}
\end{algorithm}

\section{Partial-Order-Constrained Differentiable Structure Learning (PO-DSL)}

\subsection{Partial orders as path prohibitions and quadratic overhead}
In differentiable DAG learning, $W\in\R^{d\times d}$ parameterizes a directed graph. Acyclicity is enforced
via a differentiable constraint $h(W)=0$ based on matrix power series. A general partial order $O$ can be
encoded as path prohibitions: for each $(i,j)\in O^+$, forbid any path $j\rightsquigarrow i$. A continuous
characterization yields a penalty of the form
\begin{equation}
p(W,O)=\sum_{(i,j)\in O^+}\left(\sum_{l=1}^d (W\circ W)^l\right)_{j,i},
\label{eq:path_penalty}
\end{equation}
which scales with $|O^+|$ and thus can be $\Theta(m^2)$ for a chain length $m$.

\subsection{Efficient constraint via transitive reduction and maximal paths}
Let $O^-$ be the transitive reduction and $P(O^-)$ the set of maximal paths. Define
\begin{align}
h'(W,O) &= \sum_{o\in P(O^-)} h(A(W,o)),
\label{eq:hprime_def}\\
A(W,o) &= W + \tau W_o - W\circ W_o,\quad (W_o)_{i,j}=\indic[(i,j)\in o].
\label{eq:augment}
\end{align}

\begin{theorem}[Equivalence of augmented acyclicity]
A graph $G(W)$ is a DAG and satisfies partial order constraints $O$ iff $h'(W,O)=0$.
\end{theorem}
\begin{proof}[Proof sketch]
Each maximal path $o$ is enforced by requiring acyclicity after adding $o$ as virtual edges;
maximal paths cover the transitive closure when constructed from $O^-$.
\end{proof}

\begin{algorithm}[t]
\caption{PO-DSL: Augmented acyclicity for partial-order priors~\citep{ban2024differentiable}}
\label{alg:po_dsl}
\KwIn{Adjacency $W$, partial order $O$, acyclicity function $h(\cdot)$, strength $\tau$}
$O^- \leftarrow \textsc{TransitiveReduction}(O)$\;
$P \leftarrow \textsc{MaximalPaths}(O^-)$\;
$s \leftarrow 0$\;
\For{each path $o\in P$}{
  Construct $W_o$ with $(W_o)_{i,j}=\indic[(i,j)\in o]$\;
  $A \leftarrow W + \tau W_o - W\circ W_o$\;
  $s \leftarrow s + h(A)$\;
}
\KwOut{$h'(W,O)=s$}
\end{algorithm}

\section{Unified Perspective and Hybrid Use Cases}
Our two regimes interact naturally:
\begin{itemize}
\item \textbf{From traces to priors:} inferred poset marginals from HPO-VI
  (e.g., $\hat\pi_{ij}=\Pr(i\succ j\mid D)$) can be thresholded to produce
  a partial order~$O$ for PO-DSL causal discovery.  When this prior contains
  long chains, one should enforce it via the augmented-acyclicity construction
  (Algorithm~\ref{alg:po_dsl}), not by enumerating $O^+$~\citep{ban2024differentiable}.
\item \textbf{From priors to traces:} domain ordering priors can be injected
  into HPO-VI \emph{without} enumerating $O^+$, as justified by the following theorem.
\end{itemize}

\begin{theorem}[Cover constraints suffice under dominance]
\label{thm:cover}
Let $O$ be a strict partial order on $[m]$ and let $O^-$ be its transitive reduction.
If for every $(i,j)\in O^-$ and every $k\in[K]$ we have $U_{i,k}>U_{j,k}$,
then for every $(i,j)\in O^+$ we have $i\succ_{\Dom(U)} j$.
\end{theorem}
\begin{proof}
Take $(i,j)\in O^+$.  By definition of transitive closure, there exists a
directed path $i=v_0\to v_1\to\cdots\to v_r=j$ using edges in $O^-$.
For each edge $(v_{t-1},v_t)\in O^-$, the assumption gives
$U_{v_{t-1},k}>U_{v_t,k}$ for all $k$.
By transitivity of $>$ in $\R$, we get $U_{i,k}>U_{j,k}$ for all $k$,
hence $i\succ_{\Dom(U)} j$.
\end{proof}

Theorem~\ref{thm:cover} implies that enforcing dominance only on the $|O^-|$ cover
relations (often $O(m)$ rather than $\Theta(m^2)$) automatically enforces the full
transitive closure.  In practice, one can inject an ordering prior into HPO-VI via
a soft log-prior penalty using only cover edges:
\begin{equation}
\log p(O\mid U) =
\lambda \sum_{(i,j)\in O^-}\sum_{k=1}^K
\log\sigma\!\left(\frac{U_{i,k}-U_{j,k}}{\tau}\right),
\label{eq:soft_prior}
\end{equation}
where $\sigma(\cdot)$ is the sigmoid function and $\tau>0$ is a temperature.
Alternatively, a hinge penalty $\max(0,\delta-(U_{i,k}-U_{j,k}))$ may be used.
This gives the Ban et al.~\citep{ban2024differentiable} ``avoid transitive-closure
blow-up'' philosophy \emph{inside} HPO inference, even though the latent object is $U$
rather than an adjacency $W$.

\section{Experiments (Outline)}
\textbf{Synthetic posets:} recover $h^\star$ from traces under varying $m,K,\epsilon$; compare MCMC vs.\ HPO-VI.\\
\textbf{Long-chain stress test:} demonstrate quadratic scaling of naive $|O^+|$ penalties vs.\ maximal-path scaling of augmented constraints.\\
\textbf{Real logs:} workflow/agent traces; show uncertainty-aware poset recovery and downstream execution utility.\\
\textbf{Causal discovery:} inject inferred priors into NOTEARS/DAGMA-style backbones and measure SHD/F1 improvements.

\section{Discussion}
\textbf{VI tradeoffs.}\quad
Importance-weighted objectives tighten bounds~\citep{burda2016iwae} but require
careful gradient estimation when reparameterization is unavailable.
The VIMCO-$\star$ estimator (Eq.~\ref{eq:vimco_star}) provides $\sqrt{N}$
SNR scaling~\citep{daudel2026importance}, making it the recommended default for HPO-VI.
For $\alpha>0$ (R\'{e}nyi-style bounds), generalizations of the VIMCO-$\star$ baseline
are available~\citep{daudel2026importance} but are left to future work.

\textbf{Complex partial orders.}\quad
Maximal-path factoring (Algorithm~\ref{alg:po_dsl}) remains efficient for
long chains but may grow for densely multi-chained priors.
Theorem~\ref{thm:cover} provides a complementary strategy for the HPO setting:
enforce cover constraints only, and let dominance transitivity propagate them.

\textbf{Future work.}\quad
Amortized inference over trace datasets;
soft/differentiable relaxations of the dominance map $\Dom(U)$;
joint learning of utility policies $Q$ and poset structure;
and reversible-jump extensions for inference over the latent dimension $K$.

\section{Conclusion}
We presented a unified NeurIPS-style framework connecting latent-poset trace models and partial-order-constrained differentiable graph learning. The core message is operational:
avoid enumerating transitive closure constraints for long chains by maximal-path factoring, and use
non-reparameterized importance-weighted VI when discrete poset structure breaks pathwise gradients.

\bibliographystyle{plain}
\bibliography{references}

\appendix

\section{Derivation of Score-Function Gradients for VR-IWAE}
\label{app:gradient_derivation}

We provide the full derivation of Eq.~\eqref{eq:reinforce_grad}.
Let $g(\phi)=\frac{1}{N}\sum_{n=1}^N w(z_n)^{1-\alpha}$ where $z_{1:N}\sim q_\phi$.
The objective is $\mathcal{L}^{(\alpha)}_N=\frac{1}{1-\alpha}\E[\log g(\phi)]$.

\paragraph{Step 1: Log-derivative trick.}
\begin{equation}
\nabla_\psi \mathcal{L}^{(\alpha)}_N
= \frac{1}{1-\alpha}\,\E\!\left[\frac{\nabla_\psi g(\phi)}{g(\phi)}\right].
\end{equation}

\paragraph{Step 2: Expand $\nabla_\psi g$.}
Since $z_n\sim q_\phi$ and
$w(z_n)^{1-\alpha}=\big(\frac{p(D,z_n)}{q_\phi(z_n)}\big)^{1-\alpha}$, we apply the
product rule and the score-function identity
$\nabla_\psi q_\phi(z)=q_\phi(z)\nabla_\psi\log q_\phi(z)$:
\begin{align}
\nabla_\psi g(\phi)
&= \frac{1}{N}\sum_{n=1}^N \Big[
  (1-\alpha)\,w(z_n)^{1-\alpha}\,\nabla_\psi\log w(z_n)
  + w(z_n)^{1-\alpha}\,\nabla_\psi \log q_\phi(z_n)
\Big].
\end{align}

\paragraph{Step 3: Divide by $g(\phi)$ and simplify.}
With $W_n^{(1-\alpha)}=w(z_n)^{1-\alpha}/\sum_j w(z_j)^{1-\alpha}$:
\begin{align}
\frac{\nabla_\psi g}{g}
&= (1-\alpha)\sum_n W_n^{(1-\alpha)}\nabla_\psi\log w(z_n)
  + \sum_n W_n^{(1-\alpha)}\nabla_\psi\log q_\phi(z_n).
\end{align}
Multiplying by $\frac{1}{1-\alpha}$:
\begin{align}
\nabla_\psi \mathcal{L}^{(\alpha)}_N
&= \E\!\left[
  \sum_n W_n^{(1-\alpha)}\nabla_\psi\log w(z_n)
  + \frac{1}{1-\alpha}\sum_n W_n^{(1-\alpha)}\nabla_\psi\log q_\phi(z_n)
\right].
\end{align}
Since $\sum_n W_n^{(1-\alpha)}=1$, the second term can be rewritten
using the identity $\E[\nabla_\psi\log q_\phi(z)]=0$ (applied to each sample's marginal),
yielding the REINFORCE form in Eq.~\eqref{eq:reinforce_grad} after distributing the
global $\log g$ signal across all samples.

\section{VIMCO-$\star$ Derivation and Variance Analysis}
\label{app:vimco_star}

\paragraph{Control-variate construction.}
The high-variance REINFORCE term in~\eqref{eq:reinforce_grad} multiplies
$\nabla_\psi\log q_\phi(z_n)$ by the global signal $\frac{1}{1-\alpha}\log g$.
The VIMCO approach~\citep{mnih2016vimco} replaces this global signal with a
per-sample baseline $f_{-n}$ that depends only on $z_{-n}=(z_1,\ldots,z_{n-1},z_{n+1},\ldots,z_N)$:
since $\E[\nabla_\psi\log q_\phi(z_n)\,f(z_{-n})]=0$ (the score function has
zero expectation over $z_n$ regardless of $z_{-n}$), unbiasedness is preserved.

\paragraph{VIMCO-$\star$ baseline ($\alpha=0$).}
For the IWAE case ($\alpha=0$), the optimal leave-one-out baseline analyzed by
\citet{daudel2026importance} yields the VIMCO-$\star$ estimator:
\begin{equation}
\widehat{\nabla_\psi \mathcal{L}^{(0)}_N}
= \sum_{n=1}^N\Big[
W_n\,\nabla_\psi \log w(z_n)
- \log(1-W_n)\,\nabla_\psi \log q_\phi(z_n)
\Big],
\end{equation}
where $W_n=w(z_n)/\sum_j w(z_j)$.

\begin{proposition}[SNR scaling, {\citet[Theorem~3]{daudel2026importance}}]
\label{prop:snr}
Under mild regularity conditions, the signal-to-noise ratio of the VIMCO-$\star$
estimator at $\alpha=0$ scales as $\Theta(\sqrt{N})$, in contrast to the
$O(1/\sqrt{N})$ SNR decay of naive multi-sample REINFORCE.
\end{proposition}

\paragraph{Specialization to variational parameters.}
For $\psi\subseteq\phi$ (so $\nabla_\psi\log w=-\nabla_\psi\log q$),
the estimator simplifies to
\begin{equation}
\widehat{\nabla_\psi \mathcal{L}^{(0)}_N}
= -\sum_{n=1}^N \big(W_n+\log(1-W_n)\big)\,\nabla_\psi \log q_\phi(z_n).
\end{equation}
Note that $W_n+\log(1-W_n)\le 0$ for $W_n\in(0,1)$, so the effective per-sample signal
is non-negative (encouraging $\phi$ to increase log-probability of high-weight particles).

\section{Complexity Analysis}
\label{app:complexity}

\paragraph{Per-particle cost.}
Building the dominance adjacency $h^{(n)}=\Dom(U^{(n)})$ costs $O(m^2 K)$ via
vectorized broadcasting.  For each trace of length $T$, maintaining indegree counts
and masked row-sums for frontier computation gives $O(Tm)$ per trace.
Total per gradient step: $O(N(m^2 K + \sum_i T_i\, m))$, which is practical for
$m$ up to a few hundred with careful batching.

\paragraph{PO-DSL constraint cost.}
The naive path-penalty approach (Eq.~\ref{eq:path_penalty}) costs
$O(|O^+|\cdot d^3)$ per evaluation (one matrix-power computation per pair).
The augmented-acyclicity approach (Algorithm~\ref{alg:po_dsl}) costs
$O(|P(O^-)|\cdot d^3)$, where $|P(O^-)|$ is the number of maximal paths
in the transitive reduction---typically $O(m)$ for chain-like orders
versus $\Theta(m^2)$ for $|O^+|$.

\end{document}
